\cleardoublepage
\chapternonum{缩略词表}

\begin{center}
    % uncomment line 5 if you want to set the fontsize of the table.
    % \zihao{-4}
    % comment line 8 and uncomment line 9 if you need to set the table as mid-aligment.
    % or addjust the paras as your personal need.
    \begin{longtable}{m{2cm}m{8cm}m{5cm}}
    % \begin{longtable}{m{2cm}<{\centering}m{8cm}<{\centering}m{5cm}<{\centering}} 
        \toprule
        \textbf{英文缩写}&\textbf{英文全称}&\textbf{中文全称}\\
        \midrule
        \endfirsthead
        \toprule
        \textbf{英文缩写}&\textbf{英文全称}&\textbf{中文全称}\\
        \midrule
        \endhead 
        \bottomrule
        \endfoot
        \bottomrule
        \endlastfoot
        % =================================================
        % Abbreviations table:
        %
        % Option 1: Manual input (for few abbreviations)
        %   Example:
        ZJU&Zhejiang University&浙江大学\\
        ZJU&Zhejiang University&浙江大学\\
        ZJU&Zhejiang University&浙江大学\\
        ZJU&Zhejiang University&浙江大学\\
        ZJU&Zhejiang University&浙江大学\\
        ZJU&Zhejiang University&浙江大学\\
        ZJU&Zhejiang University&浙江大学\\
        ZJU&Zhejiang University&浙江大学\\
        %
        % Option 2: Auto-read from CSV (recommended for many abbreviations)
        %   - Usage: comment lines 28-35 (Manual input abbreviations) and uncomment lines 45-52, add abbreviations in 
		%            page/graduate/abbreviations.csv
        %   - File: page/graduate/abbreviations.csv
        %   - Columns: abbrev, english, chinese
        %   - LaTeX will automatically sort by 'abbrev' and output the table.
        %
        % Example usage (auto-sorted CSV):
        % \DTLloaddb{abbr}{page/graduate/abbreviations.csv}
        % \DTLsort{abbrev}{abbr}
        % \DTLforeach*{abbr}{\abbrev=abbrev,\english=english,\chinese=chinese}{%
        % \abbrev       & \english            & \chinese
        %     \DTLiflastrow{}{
        % \\
        %     }
        % }
	\end{longtable}
\end{center}
